\documentclass[11pt,a4paper]{article}
\usepackage[paperwidth=210mm,paperheight=1800mm,margin=1.5cm]{geometry}
\usepackage{graphicx}
\usepackage{float}
\usepackage{caption}
\usepackage{subcaption}
\usepackage{amsmath}
\usepackage[hidelinks]{hyperref}
\usepackage{xcolor}

% No page numbers
\pagestyle{empty}
\makeatletter
\let\ps@plain\ps@empty
\makeatother

\setlength{\parindent}{0pt}
\setlength{\parskip}{0.5em}
\captionsetup{font=small,labelfont=bf}

% Adjust this in one place if the plots live somewhere else.
\newcommand{\plotdir}{../results}

\newcommand{\resultpair}[5]{%
\begin{figure}[H]
    \centering
    \begin{subfigure}[t]{0.48\linewidth}
        \centering
        \includegraphics[width=\linewidth]{\plotdir/#1}
        \caption{#2}
    \end{subfigure}
    \hfill
    \begin{subfigure}[t]{0.48\linewidth}
        \centering
        \includegraphics[width=\linewidth]{\plotdir/#3}
        \caption{#4}
    \end{subfigure}
    \caption{#5}
\end{figure}
}

\title{Exercise 3: Normalizing Flows for Stellar Labels and Uncertainties}
\author{}
\date{}

\begin{document}

\maketitle

This document shows the final plots for the three models I trained on the GALAH star spectra dataset. The models predict the stellar labels $T_{\mathrm{eff}}$, $\log g$, and $\mathrm{[Fe/H]}$, and they also try to predict the uncertainty of those values using negative log-likelihood training with Jammy Flows.

I tested three versions of the model:
\begin{itemize}
    \item \textbf{Diagonal Gaussian flow}: treats the uncertainty of each target separately.
    \item \textbf{Full Gaussian flow}: allows the three target uncertainties to be correlated.
    \item \textbf{Gaussianization flow}: a more flexible flow that can model non-Gaussian output distributions.
\end{itemize}

\section*{Training Curves}

\begin{figure}[H]
    \centering
    \begin{subfigure}[t]{0.32\linewidth}
        \centering
        \includegraphics[width=\linewidth]{\plotdir/diagonal_gaussian_loss.png}
        \caption{Diagonal Gaussian}
    \end{subfigure}
    \hfill
    \begin{subfigure}[t]{0.32\linewidth}
        \centering
        \includegraphics[width=\linewidth]{\plotdir/full_gaussian_loss.png}
        \caption{Full Gaussian}
    \end{subfigure}
    \hfill
    \begin{subfigure}[t]{0.32\linewidth}
        \centering
        \includegraphics[width=\linewidth]{\plotdir/full_flow_loss.png}
        \caption{Gaussianization flow}
    \end{subfigure}
    \caption{Training and validation negative log-likelihood for the three flow models.}
    \label{fig:all-losses}
\end{figure}


\section*{1. Diagonal Gaussian Flow}

\resultpair
    {diagonal_gaussian_t_eff_pred_vs_true.png}
    {Predicted versus true $T_{\mathrm{eff}}$.}
    {diagonal_gaussian_t_eff_pull.png}
    {Pull distribution for $T_{\mathrm{eff}}$.}
    {Diagonal Gaussian results for effective temperature. On the left, the predictions are compared with the true values. On the right, the pull distribution shows how reasonable the predicted uncertainties are.}

\resultpair
    {diagonal_gaussian_log_g_pred_vs_true.png}
    {Predicted versus true $\log g$.}
    {diagonal_gaussian_log_g_pull.png}
    {Pull distribution for $\log g$.}
    {Diagonal Gaussian results for surface gravity.}

\resultpair
    {diagonal_gaussian_fe_h_pred_vs_true.png}
    {Predicted versus true $\mathrm{[Fe/H]}$.}
    {diagonal_gaussian_fe_h_pull.png}
    {Pull distribution for $\mathrm{[Fe/H]}$.}
    {Diagonal Gaussian results for metallicity.}


\section*{2. Full Gaussian Flow}

\resultpair
    {full_gaussian_t_eff_pred_vs_true.png}
    {Predicted versus true $T_{\mathrm{eff}}$.}
    {full_gaussian_t_eff_pull.png}
    {Pull distribution for $T_{\mathrm{eff}}$.}
    {Full Gaussian results for effective temperature. This version can model correlations between the uncertainties of the three predicted labels.}

\resultpair
    {full_gaussian_log_g_pred_vs_true.png}
    {Predicted versus true $\log g$.}
    {full_gaussian_log_g_pull.png}
    {Pull distribution for $\log g$.}
    {Full Gaussian results for surface gravity.}

\resultpair
    {full_gaussian_fe_h_pred_vs_true.png}
    {Predicted versus true $\mathrm{[Fe/H]}$.}
    {full_gaussian_fe_h_pull.png}
    {Pull distribution for $\mathrm{[Fe/H]}$.}
    {Full Gaussian results for metallicity.}


\section*{3. Gaussianization Flow}

\resultpair
    {full_flow_t_eff_pred_vs_true.png}
    {Predicted versus true $T_{\mathrm{eff}}$.}
    {full_flow_t_eff_pull.png}
    {Pull distribution for $T_{\mathrm{eff}}$.}
    {Gaussianization-flow results for effective temperature. This is the most flexible model, because it can learn non-Gaussian prediction shapes.}

\resultpair
    {full_flow_log_g_pred_vs_true.png}
    {Predicted versus true $\log g$.}
    {full_flow_log_g_pull.png}
    {Pull distribution for $\log g$.}
    {Gaussianization-flow results for surface gravity.}

\resultpair
    {full_flow_fe_h_pred_vs_true.png}
    {Predicted versus true $\mathrm{[Fe/H]}$.}
    {full_flow_fe_h_pull.png}
    {Pull distribution for $\mathrm{[Fe/H]}$.}
    {Gaussianization-flow results for metallicity.}


\section*{4. Example Predicted PDFs from the Gaussianization Flow}

\begin{figure}[H]
    \centering
    \includegraphics[width=0.8\linewidth]{\plotdir/full_flow_pdf_example_0.png}
    \caption{Example PDF predicted by the Gaussianization flow for one test star. The three panels show the sampled one-dimensional distributions for $T_{\mathrm{eff}}$, $\log g$, and $\mathrm{[Fe/H]}$, together with the true value and the sample mean.}
    \label{fig:full-flow-pdf-0}
\end{figure}

\begin{figure}[H]
    \centering
    \includegraphics[width=0.8\linewidth]{\plotdir/full_flow_pdf_example_1.png}
    \caption{Second example of a predicted PDF from the Gaussianization flow.}
    \label{fig:full-flow-pdf-1}
\end{figure}

\begin{figure}[H]
    \centering
    \includegraphics[width=0.8\linewidth]{\plotdir/full_flow_pdf_example_2.png}
    \caption{Third example of a predicted PDF from the Gaussianization flow.}
    \label{fig:full-flow-pdf-2}
\end{figure}

\end{document}
